\subsection{整数集合中的区间}

每个整数集合，作为基数集合，都是良序的（§ 3，第 2 号，定理 1）。此外，对于每个整数 \(a\) ，关系“ \(x\) 是一个基数，且 \(x \leqslant a\) '' 在 \(x\) （第3节，第2号，定理1后的注记）中是收集性的，并且满足此关系的 \(x\) 的集合是一个整数集合（第4节，第2号，命题2），因此可以记为 {[}0, a{]}.

命题4。设 \(a\) 与 \(b\) 是整数。则映射 \(x \to a + x\) 是区间 \([0, b]\) 到区间 \([a, a + b]\) 上的严格递增同构，而 \(y \to y - a\) 是逆同构。

显然，关系 \(0 \leqslant x \leqslant b\) 蕴含 \(a \leqslant a + x \leqslant a + b\) . 由第2号命题3，映射 \(x \to a + x\) 严格递增（因而是单射）。最后，关系 \(a \leqslant y \leqslant a + b\) 蕴含 \(y = a + x\) ，其中 \(x \geqslant 0\) 且 \(a + x \leqslant a + b\) ，由此 \(x \leqslant b\) （第2号，命题3）。证明完成。

命题5. 如果 \(a\) 与 \(b\) 是满足 \(a \leqslant b\) 的整数，则区间 \([a, b]\) 是一个有限集合，其元素个数为 \(b - a + 1\) .

根据命题4，我们可以将自身限制在以下情形： \(a=0\) 。证明通过对 \(b\) 进行归纳。如果 \(b=0\) ，结果是显然的。关系 \(0 \leqslant x \leqslant b+1\) 等价于“ \(0 \leqslant x < b+1\) 或 \(x=b+1\) ''，且关系 \(0 \leqslant x < b+1\) 等价于 \(0 \leqslant x \leqslant b\) (§ 4，第 2 号，命题 2)；换言之，区间 \([0, b+1]\) 是 \([0, b]\) 与 \(\{b+1\}\) 的并集，且这两个集合不相交。根据归纳假设， \([0, b+1]\) 的元素个数等于 \((b+1)+1\) ，命题得证。

命题6. 对于每个有限集合 \(\mathbf{E}\) （它是全序的且具有 \(n\) 个元素 \((n \geqslant 1)\) ），存在唯一的从 \(\mathbf{E}\) 到区间 \([1, n]\) 上的同构.

由于 E 和 \([1, n]\) 都是良序的（§4，第 4 号，命题 3，推论 1）且具有相同数量的元素（命题 5），结论由 §2，第 5 号的定理 3 以及 §4，第 2 号的命题 2 的推论 2 得出。

